\methodname{} has several known limitations that we believe deserve
explicit acknowledgement.

\paragraph{Acc-2 / F1 are not state-of-the-art on MOSI.}
Even at the Acc-tuned operating point, \methodname{} reaches
$88.85\%$ Acc-2 / $88.83\%$ F1, which is roughly one point below the
contemporary CaMIB submission \citep{camib2026}.
The MAE-tuned configuration trades classification accuracy for a
$0.027$ MAE reduction; we expect that combining \methodname{}'s
IB-modulation pathway with CaMIB's causal regularisation could close
this gap, but we leave such cross-architecture studies to future work.

\paragraph{Routing collapses to argmax-language at inference.}
Although the soft routing weights are non-degenerate
(Figure~\ref{fig:routing}b), the deterministic
$\arg\max\boldsymbol{\ell}$ chooses language for $100\%$ of the MOSI
test set.
This means that the dynamic primary-modality selection of
\S\ref{sec:method:mselector} effectively manifests at training time
through the Gumbel sampling and at inference only through the soft
weights that scale auxiliary streams.
A harder constraint -- e.g.\ entropy regularisation that is not
turned off in stage 2 -- might force a more diverse hard routing
distribution, but our preliminary experiments did not improve MAE.

\paragraph{No state-of-the-art result on CH-SIMS~v2.}
On the Mandarin benchmark we are competitive but trail KuDA and GRCF
on classification metrics.
We attribute the gap to the much shorter aligned segments per token
($\!\le\!50$ subwords with $d_{a}\!=\!25$, $d_{v}\!=\!177$), which
under-exploits the IB-confidence pathway that benefits from longer
sequences with token-level uncertainty diversity.
Variable-length \methodname{} on the original SIMSv2 segment-level
features is on our roadmap.

\paragraph{Restricted to complete-modality evaluation.}
We do not report \emph{missing-modality} robustness in this paper.
The IB-confidence signal would naturally extend to a missing-token
mask (set $\conf=0$ where the modality is missing), but the routing-IB
target signal would have to be re-derived; we have not validated this
empirically.

\paragraph{Restricted to short utterances and English / Mandarin.}
All three benchmarks consist of $\le\!10$~s utterances delivered by
recorded speakers.
Whether \methodname{} generalises to longer dialogues, multi-speaker
settings, or other languages is currently unknown.
Multimodal humour and sarcasm benchmarks (HKT \citep{hkt2020},
MUStARD \citep{castro2019mustard}, UR-FUNNY
\citep{hasan2019urfunny}) would be natural next targets.

\paragraph{Compute footprint.}
The Optuna study used to obtain the reported numbers spans
$\sim\!200$ trials per dataset across multiple Tier configurations,
each trial training for $40$--$120$ epochs on a single A800.
Although every individual trial is cheap ($\le\!1$ GPU-hour for MOSI),
reproducing the entire study requires $\sim\!250$ GPU-hours of total
compute.

\paragraph{Legacy state-dict slots.}
The released checkpoints retain unused parameter slots
(e.g.\ \texttt{reverse\_proj}) inherited from an earlier development
branch.
These can be safely \texttt{strict=False} loaded but inflate the
checkpoint size by $\sim\!1\%$.
